\cotefile{11}
\scanpage{3}
% e01: top, lines 1–2
$\mathcal V$ catégorie (\unsure{Objts} $X,Y,\ldots$) avec produits de deux objets
(\unsure{et} un objet final)

% e02: lines 3–5, with additions and cancellation
$A:\quad X\longrightarrow A(X)$ foncteur cov. de $\mathcal V$ dans la catégorie
\inserted{$C$} des Algèbres \inserted{gr.}, bivariant, satisfaisant une condition d'échange
\inserted{\unsure{entre les} $f^*$ et $f_*$} \deleted{\unsure{condition}}
(\unsure{dont le besoin sera vu dans la suite|dont le besoin est vu dès le suivant}\ldots).

% e03: lines 6–7
On va définir une nouvelle catégorie, notée $\mathcal V^A$, et
deux foncteurs $\varphi,\psi$
\diagram{\begin{tikzcd}[column sep=large,row sep=large]
\mathcal V \arrow[r,"\varphi"] \arrow[dr,"A"'] & \mathcal V^A \arrow[r,"\psi"] & k\text{-mod gr.} \\
& C \arrow[ur,"\omega"'] &
\end{tikzcd}}
\margin{$\omega$ oubli de \unsure{la} structure d'Algèbres,}
rendant le carré commutatif. De plus,
le foncteur $\varphi$ induit une bijection (en fait,
une identité) sur les ens. d'objets on \deleted{\unsure{identifie}}
\[
\operatorname{Ob}\mathcal V^A=\operatorname{Ob}\mathcal V
\]

% e04: middle
On pose
\[
\operatorname{Hom}_{\mathcal V^A}(X,Y)=A(X\times Y)
\]
On définit le produit de composition ainsi
\[
\begin{aligned}
\operatorname{Hom}_{\mathcal V^A}(X,Y)\times\operatorname{Hom}_{\mathcal V^A}(Y,Z)&\longrightarrow\operatorname{Hom}_{\mathcal V^A}(X,Z)\\
(u,v)&\longmapsto v\circ u
\end{aligned}
\]
\[
A(X\times Y)\times A(Y\times Z)\longrightarrow A(X\times Z)
\]
\diagram{\begin{tikzcd}[column sep=small]
& X\times Y\times Z \arrow[dl,"p_{12}"'] \arrow[d,no head,"p_{13}"] \arrow[dr,"p_{23}"] &\\
X\times Y & X\times Z & Y\times Z
\end{tikzcd}}
\[
\boxed{v\circ u=p_{31*}\bigl(p_{32}^*(\unsure{v})\,p_{21}^*(u)\bigr).}
\]
\mspage{2}
